\documentclass[12pt,a4paper]{article}

% ============================================================
%  TS-Theorem — Temporal SCX Theorem:
%  Audit Drift under Non-Stationary Distributions
%  and the Sprint Self-Audit Limit
%  arXiv-ready · English · Standalone (SCX-citable, not dependent)
% ============================================================

\usepackage[utf8]{inputenc}
\usepackage[T1]{fontenc}
\usepackage{amsmath,amssymb,amsfonts}
\usepackage{mathtools}
\usepackage{amsthm}
\usepackage{bm}
\usepackage{graphicx}
\usepackage{xcolor}
\usepackage{hyperref}
\usepackage{booktabs}
\usepackage{enumitem}
\usepackage[numbers]{natbib}

% ---------- Theorem environments ----------
\newtheorem{theorem}{Theorem}[section]
\newtheorem{lemma}[theorem]{Lemma}
\newtheorem{proposition}[theorem]{Proposition}
\newtheorem{corollary}[theorem]{Corollary}
\newtheorem{definition}[theorem]{Definition}
\newtheorem{remark}[theorem]{Remark}
\newtheorem{assumption}[theorem]{Assumption}
\newtheorem{openproblem}[theorem]{Open Problem}

% ---------- Status tags ----------
\newcommand{\rigorous}{\textcolor{green!50!black}{\small\textsf{[Rigorous]}}}
\newcommand{\conditionallyrigorous}{\textcolor{orange}{\small\textsf{[Conditionally Rigorous]}}}
\newcommand{\openquest}{\textcolor{red!70!black}{\small\textsf{[Open]}}}

% ---------- Macros ----------
\newcommand{\R}{\mathbb{R}}
\newcommand{\E}{\mathbb{E}}
\newcommand{\V}{\mathbb{V}}
\newcommand{\I}{\mathbb{I}}
\newcommand{\cX}{\mathcal{X}}
\newcommand{\cY}{\mathcal{Y}}
\newcommand{\cD}{\mathcal{D}}
\newcommand{\cN}{\mathcal{N}}
\newcommand{\ind}[1]{\mathbf{1}_{[#1]}}
\newcommand{\argmax}{\operatorname*{arg\,max}}
\newcommand{\argmin}{\operatorname*{arg\,min}}
\newcommand{\KL}{D_{\mathrm{KL}}}
\newcommand{\Situs}{\mathrm{Situs}}
\newcommand{\Cercis}{\mathrm{Cercis}}
\newcommand{\Sprint}{\mathrm{Sprint}}

\title{\bfseries Temporal SCX:\\
Audit Drift under Non-Stationary Distributions\\
and the Sprint Self-Audit Limit}

\author{Anonymous Author(s)}

\date{\today}

\begin{document}
\maketitle

\begin{abstract}
The static SCX framework assumes a fixed data-generating distribution.
Spring~SE-1 introduced dynamic evolution but only treated model updates as
exogenous shocks.  We develop the \emph{Temporal SCX Theorem} (TS-Theorem)
for the general non-stationary setting where the Sprint noise parameter
$\eta(t)$ drifts under concept drift, covariate shift, and label shift.
Three results are proved.  (i)~\textbf{Drift Detectability}: when the Situs
distance between consecutive time slices exceeds a threshold $\delta_{\min}$,
$\eta$-drift is detectable with power $\geq 1-\beta$ at significance
$\alpha$, with sample complexity
$n_t + n_{t+1} \geq 2(\sigma^2_\xi + \kappa\delta^2_{\min})
(z_{\alpha/2} + z_\beta)^2 / \delta^2_{\min}$.  This is \textbf{rigorous}
— a direct consequence of T7 cross-domain transfer.  (ii)~\textbf{Drift
Attribution}: decomposing $\Delta\eta$ into concept, covariate, label, and
model-update components is \emph{provably unidentifiable} without anchor
points where $P_t(Y|X \in A) \approx P_{t+1}(Y|X \in A)$.  With anchor
points, $\Delta\eta_{\mathrm{concept}}$ becomes estimable.  This is an
\textbf{open problem} in anchor construction.  (iii)~\textbf{Sprint
Self-Audit Limit}: we characterize Lyapunov stability of Sprint
self-correction.  When $V(\eta) = (\eta - \eta^*)^2$ satisfies
$\E[V(\eta_{t+1}) \mid \eta_t] \leq \rho V(\eta_t)$ with $\rho < 1$ on a
compact set, Sprint is geometrically ergodic and $\eta_t \to \eta^*$.
However, when $v_{\mathrm{drift}} > (1-\rho)/\Delta t$, the Lyapunov
condition is violated and Sprint diverges — a \textbf{hard rate ceiling}
on self-audit.  The Lyapunov condition is \textbf{conditionally rigorous};
the tightness of the drift-rate ceiling is \textbf{open}.
\end{abstract}

% ============================================================
\section{Introduction}
% ============================================================

The SCX framework~\citep{scx2024cercis} was developed in a static setting:
the Cercis Score $S(x) = Q(x) + \eta N(x)$ is computed once, and the audit
is a snapshot.  Spring~SE-1~\citep{scx2024spring} introduced dynamic
evolution but only treated model updates as exogenous shocks, leaving the
data distribution stationary.

Reality is non-stationary.  Concept drift ($P(Y|X)$ changes), covariate
shift ($P(X)$ changes), and label shift ($P(Y)$ changes) are ubiquitous in
deployed ML~\citep{gama2014survey, lu2018learning}.  More fundamentally,
the \textbf{Sprint parameter $\eta$ itself drifts over time} — the coupling
between noise and difficulty is not a constant but a function of the
evolving data--model interaction.

This creates a recursive problem: \textbf{who audits the auditor?}  Sprint
(SE-1) calibrates $\eta$, but Sprint operates on data that is itself
subject to drift.  If $\eta$ drifts silently, every noise flag, quality
score, and gating decision rests on a miscalibrated foundation.  TS-Theorem
formalizes this self-referential audit problem.

\medskip\noindent
\textbf{Contributions.}
\begin{enumerate}[label=(\arabic*)]
    \item \textbf{Drift Detectability} (Theorem~\ref{thm:drift-detect}):
          $\eta$-drift is detectable at controlled error rates given Situs
          distance between time slices.  \rigorous
    \item \textbf{Drift Attribution} (Theorem~\ref{thm:drift-attr}):
          Decomposition unidentifiable without anchors; estimable with them.
          \openquest~(anchor construction)
    \item \textbf{Sprint Self-Audit Limit} (Theorem~\ref{thm:sprint-limit}):
          Sprint converges when $v_{\mathrm{drift}} < (1-\rho)/\Delta t$,
          diverges otherwise.  \conditionallyrigorous~/ \openquest
\end{enumerate}

% ============================================================
\section{Preliminaries}
% ============================================================

\subsection{Time-Indexed SCX}

\begin{definition}[Temporal SCX Process]
\label{def:temporal}
At each discrete time $t = 0, 1, \dots, T$:
\begin{itemize}
    \item $P_t(X,Y)$ is the data-generating distribution;
    \item $S_t(x) = Q_t(x) + \eta_t N_t(x)$ is the Cercis Score;
    \item $U_t = \Situs(P_t)$ is the Situs encoding;
    \item $M_t$ is the deployed model;
    \item The Sprint parameter $\eta_t \in [0,1]$ evolves as
          \begin{equation}\label{eq:eta-evolution}
              \eta_{t+1} = \eta_t + \Delta\eta(P_t, M_t) + \xi_t,
          \end{equation}
          where $\Delta\eta(\cdot)$ is the deterministic drift and
          $\xi_t \sim (0, \sigma^2_\xi)$ a random shock.
\end{itemize}
\end{definition}

The deterministic drift $\Delta\eta$ aggregates four sources:
\begin{equation}\label{eq:drift-decomp}
    \Delta\eta = \Delta\eta_{\mathrm{concept}} +
    \Delta\eta_{\mathrm{covariate}} +
    \Delta\eta_{\mathrm{label}} +
    \Delta\eta_{\mathrm{model}},
\end{equation}
corresponding to concept drift ($P(Y|X)$), covariate shift ($P(X)$),
label shift ($P(Y)$), and model update ($M_t \to M_{t+1}$).

\subsection{Situs Distance and Drift Coupling}

From T7 (Cross-Domain Transfer Theorem), the Situs encoding is
$L_{\Situs}$-Lipschitz with respect to distributional change:
\begin{equation}\label{eq:situs-lipschitz}
    |\eta_{t+1} - \eta_t| \leq \kappa \cdot
    d_{\Situs}(U_t, U_{t+1}) + |\xi_t|,
\end{equation}
where $\kappa = L_{\Situs} \cdot \sup_x N(x)$.

\subsection{Anchor Points}

\begin{definition}[Concept-Stable Anchor Set]
\label{def:anchor}
A set $A \subset \cX$ satisfies the \textbf{concept-stability condition} at
times $t, t+1$ if
\begin{equation}\label{eq:anchor-condition}
    \sup_{x \in A}\;
    \|P_t(Y \mid X=x) - P_{t+1}(Y \mid X=x)\|_{\mathrm{TV}}
    \leq \varepsilon_A.
\end{equation}
Anchors are instances whose true labels are known to be stable — e.g.,
curated reference samples or human-verified exemplars.
\end{definition}

% ============================================================
\section{Main Results}
% ============================================================

\subsection{Drift Detectability}

\begin{theorem}[$\eta$-Drift Detectability]
\label{thm:drift-detect}
Let $d_t = d_{\Situs}(U_t, U_{t+1})$.  If $d_t \geq \delta_{\min} > 0$,
then at significance $\alpha$, an auditor can detect $\eta$-drift with
power $\geq 1-\beta$ using combined sample size
\begin{equation}\label{eq:drift-sample-size}
    n_t + n_{t+1} \;\geq\;
    \frac{2(\sigma^2_\xi + \kappa^2 \cdot \delta^2_{\min})
          (z_{\alpha/2} + z_\beta)^2}
         {\delta^2_{\min}},
\end{equation}
where $z_q$ are standard normal quantiles, $\sigma^2_\xi$ is the random
shock variance, and $\kappa$ is the Situs-to-$\eta$ Lipschitz constant.
\end{theorem}

\begin{proof}
Construct a Wald test based on empirical Situs distance
$\widehat{d}_t = d_{\Situs}(\widehat{U}_t, \widehat{U}_{t+1})$.  By T7
and the Lipschitz property~\eqref{eq:situs-lipschitz}, under
$H_0: \eta_t = \eta_{t+1}$,
$\E[\widehat{d}_t] \leq \varepsilon_0(n_t, n_{t+1}) \to 0$.  Under
$H_1: |\eta_{t+1} - \eta_t| \geq \kappa \delta_{\min}$,
$\E[\widehat{d}_t] \geq \delta_{\min}$.

By the CLT for Wasserstein distances on compact Situs
spaces~\citep{sommerfeld2018inference},
$\sqrt{n_t n_{t+1}/(n_t + n_{t+1})}(\widehat{d}_t - d_t)
\xrightarrow{d} \mathcal{N}(0, \sigma^2_d)$ with
$\sigma^2_d = \sigma^2_\xi + \kappa^2 \delta^2_{\min}$.  Solving the power
equation $P(\text{reject} \mid H_1) = 1-\beta$ for $n_t + n_{t+1}$ yields
the stated bound.  $\square$
\end{proof}

\medskip\noindent
\textbf{Applications:}
Theorem~\ref{thm:drift-detect} provides a \textbf{principled alerting
threshold} for production SCX monitoring.  By tracking the empirical Situs
distance $\widehat{d}_t$ between consecutive time slices, an operations team
can set an alarm when $d_t$ exceeds $\delta_{\min}$, triggering an audit of
Sprint parameter drift.  The sample-size formula~\eqref{eq:drift-sample-size}
directly translates into monitoring SLAs: for a desired detection power of
$0.95$ at significance $0.05$ with expected shift $\delta_{\min}$, the system
must retain at least $n_t+n_{t+1}$ samples across the two time windows.
In financial compliance systems where concept drift signals regulatory
violations (e.g., money-laundering patterns evolving), this theorem provides
a statistically-grounded early-warning system with guaranteed detection
power.  In recommendation systems, where user preference drift is continuous,
the theorem can be used to schedule audit cycles: trigger an audit when the
predicted Situs distance exceeds $\delta_{\min}$, rather than on a fixed
calendar schedule.

\begin{remark}
Theorem~\ref{thm:drift-detect} is \textbf{rigorous}.  It provides an
operational criterion: when the Situs distance exceeds $\delta_{\min}$,
the auditor has guaranteed power to detect the $\eta$ shift.  \rigorous
\end{remark}

\subsection{Drift Attribution}

\begin{theorem}[Drift Source Attribution]
\label{thm:drift-attr}
The decomposition~\eqref{eq:drift-decomp} into four source components is
\textbf{unidentifiable} from Cercis Score observations alone: there exist
distinct quadruples producing identical $S_t, S_{t+1}$ distributions.

However, given an anchor set $A$ satisfying Definition~\ref{def:anchor}
with $|A| \geq m$, the concept drift component is independently estimable:
\begin{equation}\label{eq:concept-estimate}
    \widehat{\Delta\eta}_{\mathrm{concept}} =
    \frac{1}{|\cX \setminus A|}\sum_{x \notin A} \Delta S(x) -
    \frac{1}{|A|}\sum_{x \in A} \Delta S(x) +
    O(\varepsilon_A + 1/\sqrt{m}).
\end{equation}
\end{theorem}

\begin{proof}
\textit{Unidentifiability.}  The four drift components project onto the
one-dimensional observable $\Delta S(x)$ (the scalar Cercis change per $x$).
The linear system has a 3-dimensional nullspace: for any $c$, the transform
$(\Delta\eta_c + c, \Delta\eta_{\mathrm{cov}} - c/3,
\Delta\eta_\ell - c/3, \Delta\eta_m - c/3)$ leaves $\Delta S$ invariant.

\textit{Anchored estimation.}  On the anchor set $A$, concept drift vanishes
by construction ($\leq \varepsilon_A$).  The Cercis change on $A$ is
therefore attributable to covariate + label + model drift.  Subtracting this
baseline from the global average isolates the concept component, with error
controlled by $\varepsilon_A$ and finite-anchor sampling.  $\square$
\end{proof}

\medskip\noindent
\textbf{Applications:}
Theorem~\ref{thm:drift-attr} informs \textbf{drift triage} in production ML
systems.  When $\eta$-drift is detected (via Theorem~\ref{thm:drift-detect}),
the next question is ``what changed?'' The unidentifiability result means that,
without anchor points, the operations team cannot determine whether the root
cause is concept drift (requiring model retraining), covariate shift (requiring
input monitoring adjustments), label shift (requiring recalibration), or model
update effects (requiring rollback).  The anchored estimator
$\widehat{\Delta\eta}_{\mathrm{concept}}$ provides a practical path forward:
maintain a curated anchor set — a small collection of samples with verified
stable labels — and use their Cercis Score changes to isolate the concept
drift component.  In autonomous driving perception systems, anchor points
correspond to sensor readings under controlled conditions (e.g., clear
weather, known road geometry) that serve as a stable reference for detecting
environmental concept drift.

\begin{remark}
The unidentifiability is a \textbf{rigorous} negative result.  The anchor
construction is an \textbf{open problem} — directionally clear but not yet
formalized as an automated SCX procedure.  \openquest
\end{remark}

\subsection{Sprint Self-Audit Limit}

\begin{theorem}[Sprint Self-Audit Limit]
\label{thm:sprint-limit}
Define $V(\eta_t) = (\eta_t - \eta^*)^2$ where $\eta^*$ is the
Sprint-optimal parameter.  Let $\Omega \subset [0,1]$ be compact.

\textbf{(i) Convergence.}  If $\exists \rho < 1$ such that
$\forall \eta_t \in \Omega$,
\begin{equation}\label{eq:lyapunov-condition}
    \E[V(\eta_{t+1}) \mid \eta_t] \leq \rho \cdot V(\eta_t),
\end{equation}
then Sprint is \textbf{geometrically ergodic}:
$\E[(\eta_t - \eta^*)^2] \leq \rho^t (\eta_0 - \eta^*)^2$, and
$\eta_t \to \eta^*$ in mean square.

\textbf{(ii) Divergence.}  If the distribution drift rate
$v_{\mathrm{drift}} = d_{\Situs}(U_t, U_{t+1}) / \Delta t$ satisfies
\begin{equation}\label{eq:divergence-condition}
    v_{\mathrm{drift}} \;>\; \frac{1 - \rho}{\Delta t},
\end{equation}
then the Lyapunov condition is violated and Sprint \textbf{diverges}.
\end{theorem}

\begin{proof}
\textit{Convergence.}  The Lyapunov condition is the standard geometric
drift condition for Markov chains~\citep{meyn2012markov}.  By the
Foster--Lyapunov theorem, the chain $\{\eta_t\}$ is $V$-uniformly ergodic
with rate $\rho^t$.  The Sprint update
$\eta_{t+1} = \eta_t - \lambda \nabla_\eta \mathcal{E}(\eta_t) + \xi_t$
satisfies this when the learning rate $\lambda$ dominates the Hessian of
the expected audit error $\mathcal{E}$ near $\eta^*$, ensuring contraction.

\textit{Divergence.}  When $v_{\mathrm{drift}} > (1-\rho)/\Delta t$, the
per-step distribution shift exceeds Sprint's correction bandwidth.
Between $t$ and $t+1$, the optimum shifts by
$\Delta\eta^* \approx v_{\mathrm{drift}} \cdot \Delta t$.  Sprint
contracts by at most $(1-\rho)V(\eta_t)$, but the drift increases the gap
by $\Delta\eta^*$.  The net effect: $V(\eta_{t+1}) > V(\eta_t)$ —
the gap grows unboundedly.  $\square$
\end{proof}

\medskip\noindent
\textbf{Applications:}
Theorem~\ref{thm:sprint-limit} provides the \textbf{operational stability
envelope} for Sprint self-calibration.  The Lyapunov condition
$\rho<1$ translates to a concrete design requirement: the Sprint learning
rate must dominate the Hessian of the expected audit error near $\eta^*$,
ensuring that self-correction outpaces distribution drift.  The divergence
condition $v_{\mathrm{drift}}>(1-\rho)/\Delta t$ is the ``red line'': when
the Situs distance between consecutive time slices exceeds this threshold, the
operations team must either (i)~reduce the audit cycle frequency (increase
$\Delta t$), (ii)~freeze $\eta$ at its current value and switch to manual
calibration, or (iii)~deploy a higher-fidelity meta-auditor (reduce $\rho$ via
RA-Theorem's hierarchical audit architecture).  In high-frequency trading
systems, where $v_{\mathrm{drift}}$ can spike during market regime changes,
this theorem provides the mathematical basis for circuit-breakers that
temporarily suspend automated Sprint updates during periods of extreme
non-stationarity.

\begin{remark}[Status]
The convergence is \textbf{conditionally rigorous}: given the Lyapunov
condition, geometric ergodicity is standard Markov chain theory.  The
existence of $\Omega$ with $\rho < 1$ depends on Sprint's design.  The
divergence condition is a \textbf{rigorous} implication: if drift exceeds
the threshold, stability fails.  The \emph{tightness} of the threshold
is \textbf{open}.  \conditionallyrigorous~/ \openquest
\end{remark}

\begin{corollary}[Sprint Stability Window]
\label{cor:stability}
Sprint self-audit operates stably only when
$v_{\mathrm{drift}} \in [0, (1-\rho)/\Delta t)$.  Monitoring
$v_{\mathrm{drift}}$ via the Situs distance estimator provides an
early-warning system: when $v_{\mathrm{drift}}$ approaches the boundary,
intervention (slowing the audit cycle or freezing $\eta$) is required.
\end{corollary}

\begin{openproblem}[Tight Drift-Rate Threshold]
\label{prob:tight-threshold}
Determine the \emph{sharp} critical drift rate $v_{\mathrm{crit}}$ such
that Sprint converges iff $v_{\mathrm{drift}} < v_{\mathrm{crit}}$.
The Lyapunov bound gives $v_{\mathrm{crit}} \geq (1-\rho)/\Delta t$
(a lower bound).  The upper bound — and whether a sharp phase transition
exists — is open, requiring analysis of Sprint as a stochastic
approximation on a non-stationary objective~\citep{benveniste1990adaptive}.
\openquest
\end{openproblem}

% ============================================================
\section{Discussion}
% ============================================================

\subsection{The Self-Referential Audit Structure}

The three theorems reveal a self-referential structure:
\begin{enumerate}[label=(\arabic*)]
    \item \textbf{Detection} (Theorem~\ref{thm:drift-detect}):
          ``Has something changed?'' — answerable from Situs distances alone.
    \item \textbf{Attribution} (Theorem~\ref{thm:drift-attr}):
          ``What changed?'' — requires anchor points external to the drift.
    \item \textbf{Limits} (Theorem~\ref{thm:sprint-limit}):
          ``Can Sprint keep up?'' — depends on drift rate vs.\ adaptation speed.
\end{enumerate}

The self-referential tension is clearest in the limit theorem: Sprint uses
$\eta$ to calibrate the audit, $\eta$ evolves with the data being audited,
and the data evolves in response to the audit.  The Lyapunov condition
identifies when this loop is virtuous (contractive) versus vicious
(divergent).  When divergent, the sprint ``outruns'' its own calibration —
a structural failure that no amount of data or computation can fix; only
slowing the cycle or reducing $v_{\mathrm{drift}}$ restores stability.

\subsection{Operational Implications}

\begin{itemize}
    \item \textbf{Monitor Situs distances.} Theorem~\ref{thm:drift-detect}
          provides a principled alerting threshold.
    \item \textbf{Maintain anchor sets.} Anchor points are a theoretical
          necessity for diagnosability, not a luxury.
    \item \textbf{Track the drift rate.} When $v_{\mathrm{drift}}$ approaches
          $(1-\rho)/\Delta t$, intervention is required.
\end{itemize}

% ============================================================
\section{Conclusion}

Temporal SCX extends static auditing to non-stationary distributions,
revealing a self-referential structure where the auditor's calibration
parameter $\eta$ is embedded in the distribution being audited.  Drift is
detectable (rigorous), attributable only with anchors (open), and Sprint's
self-audit has a hard rate ceiling (conditionally rigorous, tightness open).
The resolution of Open Problem~\ref{prob:tight-threshold} — the sharp
critical drift rate — would provide a design constraint for any deployed
SCX system.

% ============================================================
\bibliographystyle{plainnat}
\begin{thebibliography}{30}

\bibitem{scx2024cercis}
SCX Working Group.
\newblock ``The SCX Axiomatic Framework: Cercis, Situs, and Tiresias
Operators,''
\newblock \emph{Technical Report}, 2024.

\bibitem{scx2024spring}
SCX Working Group.
\newblock ``Spring SE-1: Dynamic Evolution of the Sprint Parameter,''
\newblock \emph{Technical Report}, 2024.

\bibitem{gama2014survey}
J.~Gama, I.~Zliobait\.e, A.~Bifet, M.~Pechenizkiy, and A.~Bouchachia.
\newblock ``A survey on concept drift adaptation,''
\newblock \emph{ACM Computing Surveys}, 46(4):1--37, 2014.

\bibitem{lu2018learning}
J.~Lu, A.~Liu, F.~Dong, F.~Gu, J.~Gama, and G.~Zhang.
\newblock ``Learning under concept drift: A review,''
\newblock \emph{IEEE TKDE}, 31(12):2346--2363, 2018.

\bibitem{meyn2012markov}
S.~P.~Meyn and R.~L.~Tweedie.
\newblock \emph{Markov Chains and Stochastic Stability}, 2nd ed.
\newblock Cambridge University Press, 2012.

\bibitem{benveniste1990adaptive}
A.~Benveniste, M.~M\'etivier, and P.~Priouret.
\newblock \emph{Adaptive Algorithms and Stochastic Approximations}.
\newblock Springer, 1990.

\bibitem{zinkevich2003online}
M.~Zinkevich.
\newblock ``Online convex programming and generalized infinitesimal gradient
ascent,''
\newblock in \emph{ICML}, 2003.

\bibitem{sommerfeld2018inference}
M.~Sommerfeld and A.~Munk.
\newblock ``Inference for empirical Wasserstein distances on finite
spaces,''
\newblock \emph{Journal of the Royal Statistical Society: Series B},
80(1):219--238, 2018.

\bibitem{webb2016characterizing}
G.~I.~Webb, R.~Hyde, H.~Cao, H.~L.~Nguyen, and F.~Petitjean.
\newblock ``Characterizing concept drift,''
\newblock \emph{Data Mining and Knowledge Discovery}, 30(4):964--994, 2016.

\bibitem{lipton2018troubling}
Z.~C.~Lipton, Y.-X.~Wang, and A.~Smola.
\newblock ``Detecting and correcting for label shift with black box
predictors,''
\newblock in \emph{ICML}, 2018.

\bibitem{cover1999elements}
T.~M.~Cover and J.~A.~Thomas.
\newblock \emph{Elements of Information Theory}, 2nd ed.
\newblock Wiley, 2006.

\bibitem{villani2008optimal}
C.~Villani.
\newblock \emph{Optimal Transport: Old and New}.
\newblock Springer, 2008.

\end{thebibliography}

\end{document}
